\documentclass[11pt]{article}

\usepackage[utf8]{inputenc}
\usepackage[T1]{fontenc}
\usepackage{lmodern}
\usepackage{geometry}
\usepackage{amsmath,amssymb,amsfonts,amsthm,mathtools}
\usepackage{bm}
\usepackage{enumitem}
\usepackage{microtype}
\usepackage{hyperref}
\usepackage{titlesec}
\usepackage{booktabs}
\usepackage{array}
\usepackage{setspace}
\usepackage{mathrsfs}
\usepackage{physics}

\geometry{margin=1in}
\hypersetup{colorlinks=true, linkcolor=black, urlcolor=blue, citecolor=black}
\setlist[enumerate]{topsep=0.4em,itemsep=0.35em}
\setlist[itemize]{topsep=0.3em,itemsep=0.25em}
\titleformat{\section}{\large\bfseries}{\thesection.}{0.5em}{}
\titleformat{\subsection}{\normalsize\bfseries}{\thesubsection.}{0.5em}{}
\setstretch{1.06}

\newcommand{\Aether}{\AE{}ther}
\newcommand{\Aflow}{\AE{}ther-flow}
\newcommand{\TheoryName}{\Aflow{} Interpretation of Relativity}
\newcommand{\TheoryProgram}{\Aether{} / \Aflow{} framework}
\newcommand{\Sub}{\mathcal{A}}
\newcommand{\BridgeMetric}{\mathfrak{g}}
\newcommand{\Ell}{\mathcal{L}}

\newtheorem{definition}{Definition}
\newtheorem{proposition}{Proposition}
\newtheorem{remark}{Remark}
\newtheorem{corollary}{Corollary}
\newtheorem{theorem}{Theorem}

\title{\textbf{The \TheoryName{}}\\[0.4em]
\large Higher-Derivative Quartic Branch Strict-Branch Decay Test and First Nonlinear-Stability Obstruction}
\author{}
\date{}

\begin{document}

\maketitle

\begin{abstract}
The explicit minimal quartic-compatible completion \(S_{\mathrm{min},4}\) now has a first local small-data theorem, a corrected coercive energy, a \(T\sim \varepsilon^{-1}\) bootstrap on the same strict auxiliary elliptic-control branch, and a same-branch test showing that \(\kappa_\theta>0\) and \(\kappa_\Omega>0\) are structural for the present local solvability architecture. The next honest question is therefore narrower again: can that corrected strict-branch energy be paired with decay or other integrable nonlinear control strongly enough to yield a genuine small-data nonlinear stability theorem on the same branch?

This note performs that test at the same disciplined scope. It assumes the best linear-compatible decay one would naturally ask for on the current branch: wave-type \(t^{-1}\) decay for the Einstein variables and Schr\"odinger-type \(t^{-3/2}\) decay for the quartic scalar channel generated by
\[
\partial_t^2\sigma + \frac{\lambda_4}{m_S^2M_*^2}\Delta^2\sigma = 0.
\]
Under that decay package, the genuinely propagating quadratic Einstein--quartic-scalar interactions are time-integrable. In particular, the quartic scalar stress
\[
\Theta_{ij}^{(\sigma)}=
\mathcal O\!\left(\frac{\lambda_4}{M_*^2}D_iD_j\sigma\,\Delta_\gamma\sigma\right)
+ \mathcal O_{\mathrm{l.o.t.}}
\]
and the lower-order curvature corrections in the scalar equation do not create the first obstruction.

The obstruction appears earlier and is more rigid. Writing \(n\equiv N-1\), the maximal-gauge lapse equation takes the elliptic form
\[
-\Delta_\gamma n + b_N\,n = \mathcal S_N,
\]
where \(\mathcal S_N\) contains the nonnegative leading pieces
\[
\frac{c_\pi}{m_S^2}\pi_\sigma^2
+ c_\Delta\frac{\lambda_4}{M_*^2}(\Delta_\gamma\sigma)^2
+ c_K |K|^2.
\]
For decaying data with nonzero monopole
\[
\mathfrak m_N(t)\equiv \int_{\mathbb R^3}\mathcal S_N(t,y)\,dy,
\]
the elliptic inverse produces the standard Coulomb tail
\[
n(t,x)=\frac{\mathfrak m_N(t)}{4\pi|x|}+\mathcal O(|x|^{-2}).
\]
The scalar equation still contains the background-slope forcing
\[
(\partial_t-\Ell_\beta)\sigma = N\frac{\pi_\sigma}{m_S^2} + \sigma_0 n,
\]
so on the wave zone \(|x|\sim t\) one obtains a nonintegrable \(t^{-1}\) source unless a new renormalization or monopole cancellation is supplied. Therefore the present branch does \emph{not} yet yield a straightforward upgrade of the corrected coercive bootstrap to a genuine nonlinear stability theorem for decay back to the exact flat background. The first obstruction is long-range auxiliary/gauge feedback through the elliptic lapse, not the quartic principal operator itself.
\end{abstract}

\section{Introduction}

The higher-derivative quartic branch has now reached a much sharper nonlinear status than before. The active manuscript line already contains one explicit minimal quartic-compatible nonlinear completion, one reduced Einstein--quartic-scalar \(3+1\) system, one first local small-data well-posedness theorem on a strict auxiliary elliptic-control branch, one corrected coercive energy yielding a \(T\sim \varepsilon^{-1}\) bootstrap on that same branch, and one same-branch relaxation test showing that the ordered-flow coefficients \(\kappa_\theta>0\) and \(\kappa_\Omega>0\) are structurally tied to the present local solvability architecture.

That changes the meaning of the next step again. The missing question is no longer whether a branch exists, whether a local theorem exists, or whether the strict elliptic assumptions can be weakened by a small variation of the same proof. The next burden is to test whether the corrected strict-branch energy can be combined with decay or other integrable nonlinear control strongly enough to close a genuine small-data nonlinear stability theorem on the branch already recorded.

\paragraph{Framework and claim boundary.}
\TheoryProgram{} denotes the broader \Aether{} / \Aflow{} framework built on the claims that the \Aether{} is the underlying four-dimensional substrate of reality, the \Aflow{} is its intrinsic ordered motion, observed three-dimensional space is the local experiential slice of that deeper substrate, S-time is the experienced order of change arising from matter, light, and the \Aflow{}, and observed expansion is the three-dimensional appearance of deeper four-dimensional ordered motion. Within that framework, \TheoryName{} denotes the exact-closure theory in which the effective gravitational dynamics are adopted to be exactly Einsteinian with universal matter coupling. In the active sequence, \emph{adoption} means use of that established relativistic dynamics without claiming substrate derivation, while \emph{derivation} is reserved for a first-principles recovery from explicit substrate variables.

The present note stays within that boundary. It does not claim a nonlinear stability theorem. It does not claim global instability of the current branch. It performs the next narrower test only: if one supplements the corrected strict-branch energy by the natural linear-compatible decay package, does the bootstrap close? The answer at present scope is no, and the first obstruction can already be isolated.

\section{The Strict-Branch Upgrade Test}

The strict branch already fixes the reduced variables, gauge, and coercive energy architecture. The only additional issue for a genuine nonlinear stability theorem is whether the nonlinear remainder can be made time-integrable.

\begin{definition}[Strict-branch corrected energy]
Let \(\mathfrak F_N(t)\) denote the corrected coercive functional from the strict-branch bootstrap note, equivalent there to the local control energy \(\mathfrak E_N(t)\) for \(N\ge 6\):
\begin{equation}
\frac{1}{2}\mathfrak E_N(t)
\le
\mathfrak F_N(t)
\le
\frac{3}{2}\mathfrak E_N(t).
\label{eq:Fequiv}
\end{equation}
\end{definition}

\begin{definition}[Linear-compatible decay package]
Fix a strict-branch solution on \([0,T]\). The solution is said to satisfy the \emph{linear-compatible decay package} if there exists \(C_\ast>0\) such that for every \(t\in[0,T]\),
\begin{align}
\mathfrak D_g(t)
&\equiv
\|\partial(\gamma-\delta)(t)\|_{W^{2,\infty}}
+ \|K(t)\|_{W^{1,\infty}}
+ \|\beta(t)\|_{W^{2,\infty}}
\le
C_\ast \varepsilon (1+t)^{-1},
\label{eq:Dg}
\\
\mathfrak D_\sigma(t)
&\equiv
\|\pi_\sigma(t)\|_{W^{1,\infty}}
+ \|D^2\sigma(t)\|_{W^{1,\infty}}
\le
C_\ast \varepsilon (1+t)^{-3/2},
\label{eq:Ds}
\\
\mathfrak D_{\mathrm{aux}}(t)
&\equiv
\|Q(t)\|_{W^{2,\infty}}
+ \|\chi(t)\|_{W^{2,\infty}}
+ \|W^T(t)\|_{W^{2,\infty}}
\le
C_\ast \varepsilon (1+t)^{-1-\delta_\ast}
\label{eq:Daux}
\end{align}
for some \(\delta_\ast>0\).
\end{definition}

\begin{remark}
Definition 2 is deliberately a \emph{test} package, not a theorem. The wave-type \(t^{-1}\) rate in \eqref{eq:Dg} is the natural first decay one would seek for the Einstein sector. The scalar rate in \eqref{eq:Ds} is the corresponding Schr\"odinger-type decay suggested by the linear quartic scalar reduction
\[
\partial_t^2 \sigma
+ \frac{\lambda_4}{m_S^2M_\ast^2}\Delta^2\sigma
=0,
\]
whose dispersion relation is \(\omega(\xi)\simeq |\xi|^2\). The point of the present note is not to prove \eqref{eq:Dg}--\eqref{eq:Daux}; it is to see whether such decay would even be enough.
\end{remark}

\section{Propagating Quadratic Terms Are Integrable Under the Test Package}

The explicit-completion and local-theorem notes already isolate the leading nonlinear structures that can enter the corrected energy estimate:
\begin{enumerate}
    \item the quartic scalar stress in the Einstein equations,
    \item the curvature-dependent lower-order corrections in the scalar equation,
    \item transport and commutator terms involving \(\beta\), \(K\), and the auxiliary variables.
\end{enumerate}
Under Definition 2 these terms are not the first obstruction.

\begin{proposition}[Integrable propagating remainders]
Assume the strict branch, the corrected coercive equivalence \eqref{eq:Fequiv}, and the linear-compatible decay package \eqref{eq:Dg}--\eqref{eq:Daux}. Then every quadratic remainder in the commuted strict-branch energy estimate arising from:
\begin{enumerate}
    \item the quartic scalar stress
    \[
    \Theta_{ij}^{(\sigma)}
    =
    \mathcal O\!\left(\frac{\lambda_4}{M_\ast^2}D_iD_j\sigma\,\Delta_\gamma\sigma\right)
    + \mathcal O_{\mathrm{l.o.t.}},
    \]
    \item the scalar momentum density
    \[
    J_i^{(\sigma)} = -\pi_\sigma D_i\sigma + \mathcal O_{\mathrm{l.o.t.}},
    \]
    \item the scalar lower-order correction
    \[
    \mathcal N_\sigma
    =
    \mathcal O(NK\pi_\sigma)
    + \mathcal O\!\left(\frac{\lambda_4}{M_\ast^2}\mathrm{Ric}[\gamma]\cdot D^2\sigma\right)
    + \mathcal O\!\left(\frac{\lambda_4}{M_\ast^2}D\,\mathrm{Ric}[\gamma]\cdot D\sigma\right)
    + \mathcal O_{\ge 2},
    \]
    \item and the already-controlled auxiliary insertions
\end{enumerate}
is bounded by
\begin{equation}
|\mathcal R_{\mathrm{prop}}(t)|
\le
C_N\Bigl(
\mathfrak D_g(t)^2
+ \mathfrak D_g(t)\mathfrak D_\sigma(t)
+ \mathfrak D_\sigma(t)^2
+ \mathfrak D_{\mathrm{aux}}(t)\bigl(\mathfrak D_g(t)+\mathfrak D_\sigma(t)\bigr)
\Bigr)\mathfrak F_N(t).
\label{eq:Rprop}
\end{equation}
In particular, the coefficient multiplying \(\mathfrak F_N(t)\) in \eqref{eq:Rprop} lies in \(L_t^1([0,T])\).
\end{proposition}

\begin{proof}
Each propagating remainder already recorded in the explicit-completion and local-theorem notes vanishes on the flat exact-closure background and contains at least one perturbative coefficient factor together with one differentiated unknown. Under \eqref{eq:Dg}--\eqref{eq:Daux}, the metric and auxiliary factors decay at least like \((1+t)^{-1}\), while the scalar principal factors \(\pi_\sigma\) and \(D^2\sigma\) decay like \((1+t)^{-3/2}\). Therefore every product in the displayed list decays at worst like \((1+t)^{-2}\) and often faster. Commuting high-order derivatives through those terms and pairing against the corrected energy variables gives \eqref{eq:Rprop} by the same Sobolev multiplication used in the coercive-bootstrap note.
\end{proof}

\begin{corollary}[The quartic principal operator is not the first obstruction]
Under Definition 2, the explicit quartic scalar channel and the genuinely propagating Einstein--quartic-scalar couplings do not by themselves prevent an \(L_t^1\)-based upgrade of the strict-branch coercive bootstrap.
\end{corollary}

\begin{proof}
Equation \eqref{eq:Rprop} shows that those terms contribute only an integrable coefficient in the corrected energy inequality.
\end{proof}

\section{Elliptic Lapse Feedback on the Strict Branch}

The obstruction comes from a different place. The scalar equation still contains the background-slope coupling
\begin{equation}
(\partial_t-\Ell_\beta)\sigma
=
N\frac{\pi_\sigma}{m_S^2}
+ \sigma_0(N-1).
\label{eq:sigmadt}
\end{equation}
The issue is therefore the decay of the lapse perturbation \(n\equiv N-1\), not the quartic principal operator.

\begin{proposition}[Strict-branch lapse equation for the perturbation]
On the strict auxiliary elliptic-control branch, the lapse perturbation \(n=N-1\) satisfies an elliptic equation of the schematic form
\begin{equation}
-\Delta_\gamma n + b_N[\gamma,K,\sigma,\pi_\sigma,\psi,\pi_\psi]\,n
=
\mathcal S_N[\gamma,K,\sigma,\pi_\sigma,Q,\chi,W^T,\psi,\pi_\psi],
\label{eq:nell}
\end{equation}
where \(b_N\ge 0\) in the flat-background small-data regime and the source has the expansion
\begin{align}
\mathcal S_N
=\;&
c_K |K|^2
+ \frac{c_\pi}{m_S^2}\pi_\sigma^2
+ c_\Delta\frac{\lambda_4}{M_\ast^2}(\Delta_\gamma\sigma)^2
+ \mathcal S_N^{\mathrm{div}}
+ \mathcal S_N^{\mathrm{int}},
\label{eq:SN}
\end{align}
with constants \(c_K,c_\pi,c_\Delta>0\), a divergence-form remainder \(\mathcal S_N^{\mathrm{div}}=D_i\mathcal J_N^i\), and an integrable remainder \(\mathcal S_N^{\mathrm{int}}\) that is at least cubic in the perturbation size.
\end{proposition}

\begin{proof}
The local well-posedness note already records that the lapse is determined slice by slice by the maximal-gauge elliptic equation. Subtracting the flat exact-closure background \(N=1\) rewrites that equation in the perturbative form \eqref{eq:nell}. In maximal slicing, the principal source terms are the Einstein quadratic geometry density together with the matter and scalar energy density. The explicit-completion note identifies the scalar energy density pieces
\[
\rho_\sigma
=
\frac{1}{2m_S^2}\pi_\sigma^2
+ \frac{\lambda_4}{2M_\ast^2}(\Delta_\gamma\sigma)^2
+ \mathcal O_{\mathrm{l.o.t.}},
\]
and those contributions enter the lapse source with positive sign in the standard maximal-gauge equation. The remaining lower-order terms either have divergence form or contain at least one extra perturbative factor. Their exact coefficients are unnecessary for the present argument.
\end{proof}

\begin{definition}[Lapse monopole]
For a strict-branch slice define the \emph{lapse monopole}
\begin{equation}
\mathfrak m_N(t)
\equiv
\int_{\mathbb R^3}\mathcal S_N(t,y)\,dy.
\label{eq:mN}
\end{equation}
\end{definition}

\begin{remark}
The current strict-branch identities do not force \(\mathfrak m_N(t)\) to vanish. Because \(\mathcal S_N\) contains the nonnegative leading pieces shown in \eqref{eq:SN}, generic small but nontrivial branch data have \(\mathfrak m_N(t)\sim \varepsilon^2\).
\end{remark}

\begin{theorem}[Generic Coulomb tail of the strict-branch lapse]
Assume the strict branch and let \(n=N-1\) solve \eqref{eq:nell} with decaying boundary condition at spatial infinity. If \(\mathfrak m_N(t)\neq 0\), then for each fixed \(t\),
\begin{equation}
n(t,x)
=
\frac{\mathfrak m_N(t)}{4\pi |x|}
+ \mathcal O(|x|^{-2})
\qquad
\text{as } |x|\to\infty.
\label{eq:coulomb}
\end{equation}
Consequently, on the wave zone \(|x|\simeq 1+t\),
\begin{equation}
|n(t,x)|
\gtrsim
|\mathfrak m_N(t)|(1+t)^{-1}
\label{eq:wavezone}
\end{equation}
unless a nongeneric cancellation forces the monopole to vanish.
\end{theorem}

\begin{proof}
In the strict branch the spatial metric stays uniformly elliptic and \(H^{N-1}\)-close to the Euclidean metric. Therefore the Green kernel of the operator \(-\Delta_\gamma+b_N\) has the same leading \(1/(4\pi|x|)\) asymptotics as the Euclidean Laplacian, with only lower-order corrections. Decompose \(\mathcal S_N\) into its integral \(\mathfrak m_N(t)\), the divergence term, and the higher-order remainder. The divergence term contributes only \(\mathcal O(|x|^{-2})\) to the far field, and the same is true for \(\mathcal S_N^{\mathrm{int}}\) because it is more localized and has one extra small factor. The monopole part therefore controls the leading tail, giving \eqref{eq:coulomb}. Evaluating at \(|x|\simeq 1+t\) yields \eqref{eq:wavezone}.
\end{proof}

\section{First Obstruction to the Stability Upgrade}

The previous section identifies the only part of the current strict branch that does not automatically improve under the linear-compatible decay test. That is enough to block a straightforward nonlinear stability theorem aimed at decay back to the exact flat background.

\begin{theorem}[Failure of the straightforward integrable upgrade test]
Assume a strict-branch solution satisfies the corrected coercive bootstrap and the linear-compatible decay package of Definition 2. Assume also that the lapse monopole is generically nonzero, so that \eqref{eq:wavezone} holds with \(|\mathfrak m_N(t)|\gtrsim c_0\varepsilon^2\) on a nontrivial time interval. Then the scalar forcing term \(\sigma_0(N-1)\) in \eqref{eq:sigmadt} is not time-integrable on the wave zone:
\begin{equation}
\int_0^T
\sup_{|x|\simeq 1+t}
|\sigma_0(N-1)(t,x)|\,dt
\gtrsim
|\sigma_0|\,c_0\varepsilon^2
\int_0^T (1+t)^{-1}\,dt.
\label{eq:logdiv}
\end{equation}
Therefore the current strict-branch package does \emph{not} close a straightforward \(L_t^1\)-based upgrade of the \(T\sim \varepsilon^{-1}\) coercive bootstrap to a genuine nonlinear stability theorem for decay back to the exact flat background.
\end{theorem}

\begin{proof}
Proposition 1 and Corollary 1 show that every genuinely propagating quadratic remainder is \(L_t^1\) under Definition 2. The only remaining nonintegrable candidate in the scalar evolution is the background-slope forcing \(\sigma_0(N-1)\). By Theorem 1, the lapse perturbation on the wave zone has size at least \(c_0\varepsilon^2(1+t)^{-1}\) whenever the monopole does not vanish. Multiplying by \(|\sigma_0|\) and integrating in time gives \eqref{eq:logdiv}. Hence the current branch does not furnish the integrable forcing needed for a direct global decay bootstrap around the exact flat background in the unrenormalized variable \(\sigma\).
\end{proof}

\begin{corollary}[First obstruction on the present branch]
At current disciplined scope, the first obstruction to promoting the strict-branch corrected coercive bootstrap to a genuine nonlinear stability theorem is long-range auxiliary/gauge feedback through the elliptic lapse equation. More precisely, the obstruction is the Coulombic tail of \(N-1\) generated by the nonzero lapse monopole and inserted into the scalar equation through the background-slope coupling \(\sigma_0(N-1)\).
\end{corollary}

\begin{proof}
This is exactly the content of Theorem 2 together with the integrable remainder statement of Proposition 1.
\end{proof}

\begin{remark}
Corollary 2 does \emph{not} say that the current branch is globally unstable in principle. It says only that the \emph{straightforward} strict-branch upgrade test fails before one reaches a full theorem. A successful future route would need at least one new ingredient: for example a renormalized scalar unknown that removes the long-range lapse phase, a gauge in which the lapse obeys a different asymptotic equation, or a different completion or auxiliary architecture that changes the source of the elliptic tail itself.
\end{remark}

\section{What This Step Establishes}

The present manuscript establishes the following.
\begin{enumerate}
    \item It tests the strict-branch corrected coercive energy against the natural linear-compatible decay package one would first try for a genuine nonlinear stability theorem.
    \item It shows that the quartic scalar principal operator and the genuinely propagating quadratic Einstein--quartic-scalar couplings are not the first obstruction under that test.
    \item It isolates the first obstruction as the long-range Coulomb tail of the elliptically determined lapse perturbation.
    \item It identifies the specific mechanism by which that tail re-enters the evolution: the background-slope forcing term \(\sigma_0(N-1)\) in the scalar equation.
\end{enumerate}

It does not establish the following.
\begin{enumerate}
    \item It does not prove that the present branch can never satisfy any nonlinear stability theorem after further renormalization or reformulation.
    \item It does not choose a new completion, a new gauge, or a new auxiliary architecture.
    \item It does not reopen the branch-selection problem.
    \item It does not alter the earlier exact-closure, local-theorem, coercive-bootstrap, or strict-elliptic-relaxation results already on record.
\end{enumerate}

\section{Conclusion}

The strict-branch coercive bootstrap has now been tested against the first natural stability-upgrade mechanism rather than merely named as a future task. Under the linear-compatible decay rates one would normally seek on this branch, the propagating quadratic Einstein--quartic-scalar interactions are already integrable. The first obstruction comes instead from the slice-wise elliptic gauge sector: the lapse perturbation acquires a generic Coulomb tail, and the scalar equation feeds that tail directly back through the background-slope term \(\sigma_0(N-1)\).

That conclusion sharpens the live program materially. The next burden is no longer to ask vaguely whether ``more decay'' might help. The first obstruction is now identified. Only after that identification does it become mathematically disciplined to consider a renormalized scalar variable, a different gauge, or a different completion or auxiliary architecture if one wants a genuine theorem beyond the current strict branch.

\end{document}
